\title{User-defined Types and Operations in \GROOVE}
\author{Arend Rensink\orcidID{0000-0002-1714-6319}}
\institute{University of Twente, \href{mailto:arend.rensink@utwente.nl}{\email{arend.rensink@utwente.nl}}}
\maketitle

\begin{abstract}
The graph transformation tool \GROOVE supports data manipulation through a signature with fixed sorts \intS, \realS, \boolS and \stringS and a collection of pre-defined operators. Graphs can contain algebra values as special nodes; attributes correspond to edges to such value nodes.

\smallskip
This setup has limitations. For instance, it is awkward to model more advanced data calculations --- for instance, advanced string or statistical operations. Moreover, the algebraic underpinning means that non-determinism (for instance, randomness) cannot be modelled. Finally, there is no support for structured data types, such as records.

\smallskip
In this paper, we describe an extension of \GROOVE that allows \emph{user-defined types and operations}, based on programmable Java classes and methods. Essentially, users can annotate their own classes and methods and make these available to \GROOVE at runtime, after which they can be employed in rule systems. Moreover, methods can be declared as \emph{indeterminate}, meaning that their outcome is not completely fixed by their arguments; this allows randomness and other forms of non-determinism.
\end{abstract}
